\chapter{2004年广东卷}
\section{单选题}
\begin{problem}
已知平面向量 \(\bs{a} = (3, 1)\)，\(\bs{b} = (x, -3)\)，且 \(\bs{a} \perp \bs{b}\)，则 \(x =\)（\quad）

\choices
    {\(-3\)}
    {\(-1\)}
    {\(1\)}
    {\(3\)}
\end{problem}

\begin{answer}
C
\end{answer}

\begin{solution}
由 \(\bs{a} \perp \bs{b}\)，得 \(\bs{a}\cdot\bs{b}=0\)，即 \(3x+1\cdot(-3)=0\). 解得 \(x=1\)，故选 C.
\end{solution}

\begin{problem}
已知 \(A = \{x \mid |2x + 1| > 3\}\)，\(B = \{x \mid x^2 + x \leqslant 6\}\)，则 \(A \cap B =\)（\quad）

\choices
    {\([-3, -2) \cup (1, 2]\)}
    {\([-3, -2) \cup (1, +\infty)\)}
    {\((-3, -2] \cup [1, 2)\)}
    {\((-\infty, -3] \cup (1, 2]\)}
\end{problem}

\begin{answer}
A
\end{answer}

\begin{solution}
由 \(|2x+1|>3\)，得 \(2x+1>3\) 或 \(2x+1<-3\)，所以 \(x>1\) 或 \(x<-2\)，即 \(A=(-\infty,-2)\cup(1,+\infty)\). 又由 \(x^2+x\leqslant6\)，得 \((x+3)(x-2)\leqslant0\)，所以 \(B=[-3,2]\). 因此 \(A\cap B=[-3,-2)\cup(1,2]\)，故选 A.
\end{solution}

\begin{problem}
设函数 \( f(x) = \begin{cases}
\frac{3x + 2}{x^2 - 4} - \frac{2}{x - 2}, & x > 2,\\
a, & x \leqslant 2,
\end{cases} \) 在 \( x = 2 \) 处连续，则 \( a =\)（\quad）

\choices
    {\(-\frac{1}{2}\)}
    {\(-\frac{1}{4}\)}
    {\(\frac{1}{4}\)}
    {\(\frac{1}{3}\)}
\end{problem}

\begin{answer}
C
\end{answer}

\begin{solution}
当 \(x>2\) 时，\(f(x)=\dfrac{3x+2}{(x-2)(x+2)}-\dfrac{2(x+2)}{(x-2)(x+2)}=\dfrac{1}{x+2}\). 因而 \(\displaystyle\lim_{x\to2^+}f(x)=\dfrac14\). 函数在 \(x=2\) 处连续，故 \(a=f(2)=\dfrac14\)，所以选 C.
\end{solution}

\begin{problem}
\(\displaystyle\lim_{n\to \infty}\left(\frac{1}{n + 1} -\frac{2}{n + 1} +\frac{3}{n + 1} -\dots +\frac{2n - 1}{n + 1} -\frac{2n}{n + 1}\right)\) 的值为（\quad）

\choices
    {\(-1\)}
    {\(0\)}
    {\(\frac{1}{2}\)}
    {\(1\)}
\end{problem}

\begin{answer}
A
\end{answer}

\begin{solution}
括号内的分子可按相邻两项分组，得到 \((1-2)+(3-4)+\cdots+((2n-1)-2n)=-n\). 因此原式为 \(\displaystyle\lim_{n\to\infty}\frac{-n}{n+1}=-1\)，故选 A.
\end{solution}

\begin{problem}
函数 \( f(x) = \sin^2\left(x + \frac{\pi}{4}\right) - \sin^2\left(x - \frac{\pi}{4}\right) \) 是（\quad）

\choices
    {周期为 \(\pi\) 的偶函数}
    {周期为 \(\pi\) 的奇函数}
    {周期为 \(2\pi\) 的偶函数}
    {周期为 \(2\pi\) 的奇函数}
\end{problem}

\begin{answer}
B
\end{answer}

\begin{solution}
利用恒等式 \(\sin^2u-\sin^2v=\sin(u+v)\sin(u-v)\)，取 \(u=x+\dfrac\pi4\)，\(v=x-\dfrac\pi4\)，得 \(f(x)=\sin2x\). 函数 \(\sin2x\) 的最小正周期为 \(\pi\)，且 \(f(-x)=-f(x)\)，所以它是周期为 \(\pi\) 的奇函数，故选 B.
\end{solution}

\begin{problem}
一台 \(X\) 型号自动机床在一小时内不需要工人照看的概率为 \(0.8000\)，有四台这种型号的自动机床各自独立工作，则在一小时内至多 \(2\) 台机床需要工人照看的概率是（\quad）

\choices
    {\(0.1536\)}
    {\(0.1808\)}
    {\(0.5632\)}
    {\(0.9728\)}
\end{problem}

\begin{answer}
D
\end{answer}

\begin{solution}
设一小时内需要工人照看的机床数为随机变量 \(X\). 单台机床需要照看的概率为 \(1-0.8000=0.2\)，所以 \(X\sim B(4,0.2)\). 因而
\(P(X\leqslant2)=P(X=0)+P(X=1)+P(X=2)=0.8^4+\mathrm C_4^1(0.2)(0.8)^3+\mathrm C_4^2(0.2)^2(0.8)^2=0.4096+0.4096+0.1536=0.9728\)，故选 D.
\end{solution}

\begin{problem}
在棱长为 \(1\) 的正方体上，分别用过共顶点的三条棱中点的平面截该正方体，则截去 \(8\) 个三棱锥后，剩下的凸多面体的体积是（\quad）

\choices
    {\(\frac{2}{3}\)}
    {\(\frac{7}{6}\)}
    {\(\frac{4}{5}\)}
    {\(\frac{5}{6}\)}
\end{problem}

\begin{answer}
D
\end{answer}

\begin{solution}
每个被截去的三棱锥的三条互相垂直的棱长均为 \(\dfrac12\)，故其体积为 \(\dfrac16\left(\dfrac12\right)^3=\dfrac1{48}\). 共截去 \(8\) 个三棱锥，截去的体积为 \(\dfrac16\)，所以剩余凸多面体的体积为 \(1-\dfrac16=\dfrac56\)，故选 D.
\end{solution}

\begin{problem}
若双曲线 \(2x^{2} - y^{2} = k\)（\(k > 0\)）的焦点到它相对应的准线的距离是 \(2\)，则 \(k =\)（\quad）

\choices
    {\(6\)}
    {\(8\)}
    {\(1\)}
    {\(4\)}
\end{problem}

\begin{answer}
A
\end{answer}

\begin{solution}
将双曲线化为标准方程 \(\dfrac{x^2}{k/2}-\dfrac{y^2}{k}=1\)，故 \(a^2=\dfrac{k}{2}\)，\(b^2=k\)，\(c^2=a^2+b^2=\dfrac{3k}{2}\)，离心率 \(e=\dfrac ca=\sqrt3\). 右焦点到相应准线的距离为 \(c-\dfrac ae=c-\dfrac{a^2}{c}=\dfrac{b^2}{c}=\sqrt{\dfrac{2k}{3}}\). 由题意 \(\sqrt{\dfrac{2k}{3}}=2\)，得 \(k=6\)，故选 A.
\end{solution}

\begin{problem}
当 \(0 < x < \frac{\pi}{4}\) 时，函数 \(f(x) = \frac{\cos^2 x}{\cos x \cdot \sin x - \sin^2 x}\) 的最小值是（\quad）

\choices
    {\(4\)}
    {\(\frac{1}{2}\)}
    {\(2\)}
    {\(\frac{1}{4}\)}
\end{problem}

\begin{answer}
A
\end{answer}

\begin{solution}
令 \(t=\tan x\)，则 \(0<t<1\)，并且
\(f(x)=\dfrac{\cos^2x}{\cos^2x(\tan x-\tan^2x)}=\dfrac{1}{t(1-t)}\). 由于 \((t-\dfrac12)^2\geqslant0\)，有 \(t(1-t)\leqslant\dfrac14\)，且分母为正，所以 \(f(x)\geqslant4\). 当 \(t=\dfrac12\)，即 \(x=\arctan\dfrac12\) 时取等号，故最小值为 \(4\)，选 A.
\end{solution}

\begin{problem}
变量 \(x\)，\(y\) 满足下列条件： \(\begin{cases}
2x + y \geqslant 12,\\
2x + 9y \geqslant 36,\\
2x + 3y = 24,\\
x \geqslant 0,\ y \geqslant 0,
\end{cases}\) 则使 \(z = 3x + 2y\) 的值最小的 \((x, y)\) 是（\quad）

\choices
    {\((4.5,3)\)}
    {\((3,6)\)}
    {\((9,2)\)}
    {\((6,4)\)}
\end{problem}

\begin{answer}
B
\end{answer}

\begin{solution}
由 \(2x+3y=24\)，得 \(y=8-\dfrac23x\). 代入不等式 \(2x+y\geqslant12\)，得 \(x\geqslant3\)；代入不等式 \(2x+9y\geqslant36\)，得 \(x\leqslant9\). 再结合 \(x\geqslant0\)，\(y\geqslant0\)，可知 \(3\leqslant x\leqslant9\). 此时 \(z=3x+2y=3x+2\left(8-\dfrac23x\right)=\dfrac53x+16\)，随 \(x\) 增大而增大，故在 \(x=3\) 时取最小值，此时 \(y=6\)，所以选 B.
\end{solution}

\begin{problem}
若 \(f(x) = \tan \left(x + \frac{\pi}{4}\right)\)，则（\quad）

\choices
    {\(f(-1) > f(0) > f(1)\)}
    {\(f(0) > f(1) > f(-1)\)}
    {\(f(1) > f(0) > f(-1)\)}
    {\(f(0) > f(-1) > f(1)\)}
\end{problem}

\begin{answer}
D
\end{answer}

\begin{solution}
因为 \(0<1-\dfrac\pi4<\dfrac\pi4\)，所以
\(f(-1)=\tan\left(-\left(1-\dfrac\pi4\right)\right)\) 介于 \(-1\) 与 \(0\) 之间，且
\(f(1)=\tan\left(\dfrac\pi2+1-\dfrac\pi4\right)=-\cot\left(1-\dfrac\pi4\right)<-1\). 另外 \(f(0)=\tan\dfrac\pi4=1\)，故 \(f(0)>f(-1)>f(1)\)，所以选 D.
\end{solution}

\begin{problem}
如图，定圆半径为 \(a\)，圆心为 \((b, c)\)，则直线 \(ax + by + c = 0\) 与直线 \(x - y + 1 = 0\) 的交点在（\quad）
\FigureLayoutDeclare{img/2004/guangdong/q12_fig1.png}{5.5cm}{25bp 16bp 60bp 4bp}
\bitmapfigure{img/2004/guangdong/q12_fig1.png}

\choices
    {第四象限}
    {第三象限}
    {第二象限}
    {第一象限}
\end{problem}

\begin{answer}
B
\end{answer}

\begin{solution}
由图可知圆在 \(y\) 轴左侧，且圆心在 \(x\) 轴上方、圆与 \(x\) 轴相交，因此 \(b+a<0\) 且 \(0<c<a\). 联立两直线，取 \(y=x+1\) 代入 \(ax+by+c=0\)，得
\(x=-\dfrac{b+c}{a+b}\)，\(y=\dfrac{a-c}{a+b}\). 因为 \(a+b<0\)，又因 \(b<-a<-c\) 从而 \(b+c<0\)，所以 \(x<0\)；同时 \(a-c>0\)，所以 \(y<0\). 交点在第三象限，故选 B.
\end{solution}

\section{填空题}
\begin{problem}
某班委会由 \(4\) 名男生与 \(3\) 名女生组成，现从中选出 \(2\) 人担任正副班长，其中至少有 \(1\) 名女生当选的概率是\fillinblank{}.（用分数作答）
\end{problem}

\begin{answer}
\(\dfrac{5}{7}\)
\end{answer}

\begin{solution}
从 \(7\) 人中选出 \(2\) 人的所有选法数为 \(\mathrm C_7^2\)，两名男生同时当选的选法数为 \(\mathrm C_4^2\). 因此至少有 \(1\) 名女生当选的概率为
\(1-\frac{\mathrm C_4^2}{\mathrm C_7^2}=1-\frac{6}{21}=\dfrac{5}{7}\)，
因此所求概率为 \(\dfrac{5}{7}\).
\end{solution}

\begin{problem}
已知复数 \(z\) 与 \((z + 2)^{2} - 8\i\) 均是纯虚数，则 \(z =\)\fillinblank{}.
\end{problem}

\begin{answer}
\(-2\i\)
\end{answer}

\begin{solution}
设 \(z=u+v\i\)，其中 \(u\)，\(v\in\R\). 因为 \(z\) 是纯虚数，所以 \(u=0\)，从而 \(z=v\i\)，且 \(v\ne0\). 于是
\((z+2)^2-8\i=(2+v\i)^2-8\i=(4-v^2)+(4v-8)\i\).
这个复数是纯虚数，故其实部为零，即
\(4-v^2=0\).
因此 \(v=2\) 或 \(v=-2\). 当 \(v=2\) 时，\((z+2)^2-8\i=0\)，不是纯虚数，舍去；当 \(v=-2\) 时，\((z+2)^2-8\i=-16\i\)，符合条件. 因此
\(z=-2\i\).
\end{solution}

\begin{problem}
由图1有面积关系：\(\frac{S_{\triangle P'A'B'}}{S_{\triangle PAB}} = \frac{PA' \cdot PB'}{PA \cdot PB}\)，则由图2有体积关系：\(\frac{V_{P-A'B'C'}}{V_{P-ABC}} =\)\fillinblank{}.

\bitmapfigure[max width=0.78\linewidth]{img/2004/guangdong/q15_fig1.png}
\end{problem}

\begin{answer}
\(\dfrac{PA' \cdot PB' \cdot PC'}{PA \cdot PB \cdot PC}\)
\end{answer}

\begin{solution}
设 \(\overrightarrow{PA'}=u\overrightarrow{PA}\)，\(\overrightarrow{PB'}=v\overrightarrow{PB}\)，\(\overrightarrow{PC'}=w\overrightarrow{PC}\)，其中 \(u\)，\(v\)，\(w>0\). 由四面体体积公式，
\[
V_{P-ABC}=\frac{1}{6}\left|\det\left(\overrightarrow{PA},\overrightarrow{PB},\overrightarrow{PC}\right)\right|
\]
将 \(\overrightarrow{PA'}\)，\(\overrightarrow{PB'}\)，\(\overrightarrow{PC'}\) 分别代入同一体积公式，行列式的三列分别乘以 \(u\)，\(v\)，\(w\)，从而
\[
\frac{V_{P-A'B'C'}}{V_{P-ABC}}=uvw=\frac{PA'}{PA}\cdot\frac{PB'}{PB}\cdot\frac{PC'}{PC}=\frac{PA'\cdot PB'\cdot PC'}{PA\cdot PB\cdot PC}
\]
因此填空为 \(\dfrac{PA'\cdot PB'\cdot PC'}{PA\cdot PB\cdot PC}\).
\end{solution}

\begin{problem}
函数 \( f(x) = \ln (\sqrt{x + 1} - 1) \)（\(x > 0\)）的反函数 \( f^{-1}(x) = \)\fillinblank{}.
\end{problem}

\begin{answer}
\(\e^{2x}+2\e^x\)
\end{answer}

\begin{solution}
令 \(y=\ln(\sqrt{x+1}-1)\). 因为 \(x>0\)，所以 \(\sqrt{x+1}-1>0\)，可得 \(\e^y=\sqrt{x+1}-1\)，即 \(\sqrt{x+1}=\e^y+1\). 从而
\(x=(\e^y+1)^2-1=\e^{2y}+2\e^y\).
交换 \(x\)，\(y\)，得到 \(f^{-1}(x)=\e^{2x}+2\e^x\). 由于原函数的值域为 \(\R\)，该反函数的定义域为 \(\R\).
\end{solution}

\section{解答题}
\begin{problem}
已知 \(\alpha\)，\(\beta\)，\(\gamma\) 成公比为 \(2\) 的等比数列（\(\alpha \in [0, 2\pi]\)），且 \(\sin \alpha\)，\(\sin \beta\)，\(\sin \gamma\) 也成等比数列. 求 \(\alpha\)，\(\beta\)，\(\gamma\) 的值.
\end{problem}

\begin{solution}
由 \(\alpha\)，\(\beta\)，\(\gamma\) 成公比为 \(2\) 的等比数列，得 \(\beta=2\alpha\)，\(\gamma=4\alpha\). 由于 \(\sin\alpha\)，\(\sin2\alpha\)，\(\sin4\alpha\) 成等比数列，按等比数列定义各项非零，故
\[
\sin^2 2\alpha=\sin\alpha\sin4\alpha=2\sin\alpha\sin2\alpha\cos2\alpha
\]
因 \(\sin\alpha\ne0\)，\(\sin2\alpha\ne0\)，先约去 \(\sin2\alpha\)，再利用 \(\sin2\alpha=2\sin\alpha\cos\alpha\)，得
\[
2\sin\alpha\cos\alpha=2\sin\alpha\cos2\alpha
\]
再约去 \(2\sin\alpha\)，得到 \(\cos\alpha=\cos2\alpha\). 因而
\(2\cos^2\alpha-1=\cos\alpha\)，\(\Longleftrightarrow (2\cos\alpha+1)(\cos\alpha-1)=0\)
若 \(\cos\alpha=1\)，则 \(\alpha=0\) 或 \(\alpha=2\pi\)，均使 \(\sin\alpha=0\)，不符合等比数列各项非零的条件，舍去. 所以 \(\cos\alpha=-\frac12\)，结合 \(\alpha\in[0,2\pi]\)，得 \(\alpha=\frac{2\pi}{3}\) 或 \(\alpha=\frac{4\pi}{3}\).

当 \(\alpha=\frac{2\pi}{3}\) 时，\((\sin\alpha,\sin\beta,\sin\gamma)=\left(\frac{\sqrt3}{2},-\frac{\sqrt3}{2},\frac{\sqrt3}{2}\right)\)，公比为 \(-1\)；当 \(\alpha=\frac{4\pi}{3}\) 时，\((\sin\alpha,\sin\beta,\sin\gamma)=\left(-\frac{\sqrt3}{2},\frac{\sqrt3}{2},-\frac{\sqrt3}{2}\right)\)，公比也为 \(-1\). 故
\(\left(\alpha,\beta,\gamma\right)=\left(\frac{2\pi}{3},\frac{4\pi}{3},\frac{8\pi}{3}\right)\)，或 \(\left(\alpha,\beta,\gamma\right)=\left(\frac{4\pi}{3},\frac{8\pi}{3},\frac{16\pi}{3}\right)\).
\end{solution}

\begin{problem}
如图，在长方体 \(ABCD - A_{1}B_{1}C_{1}D_{1}\) 中，已知 \(AB = 4\)，\(AD = 3\)，\(AA_{1} = 2\). \(E\)，\(F\) 分别是线段 \(AB\)，\(BC\) 上的点，且 \(EB = FB = 1\).

\begin{enumerate}
\item 求二面角 \(C - DE - C_1\) 的正切值；
\item 求直线 \(EC_1\) 与 \(FD_1\) 所成的余弦值.
\end{enumerate}

\bitmapfigure[max width=0.44\linewidth]{img/2004/guangdong/q18_fig1.png}
\end{problem}

\begin{solution}
以 \(A\) 为原点，\(\overrightarrow{AB}\)，\(\overrightarrow{AD}\)，\(\overrightarrow{AA_1}\) 分别为 \(x\)，\(y\)，\(z\) 轴正向建立空间直角坐标系，则
\(D=(0,3,0)\)，\(D_1=(0,3,2)\)，\(E=(3,0,0)\)，\(F=(4,1,0)\)，\(C_1=(4,3,2)\).

\begin{enumerate}
\item 平面 \(CDE\) 与平面 \(C_1DE\) 的法向量可分别取
\[
\bs{n}_1=\overrightarrow{DE}\times\overrightarrow{DC}=(0,0,12)
\]
以及
\[
\bs{n}_2=\overrightarrow{DE}\times\overrightarrow{DC_1}=(-6,-6,12)
\]
设二面角 \(C-DE-C_1\) 的平面角为 \(\theta\)，则
\[
\cos\theta=\frac{|\bs{n}_1\cdot\bs{n}_2|}{|\bs{n}_1||\bs{n}_2|}=\frac{2}{\sqrt6}
\]
所以 \(\tan\theta=\dfrac{\sqrt2}{2}\).

\item \(\overrightarrow{EC_1}=(1,3,2)\)，\(\overrightarrow{FD_1}=(-4,2,2)\). 设两直线所成的角为 \(\varphi\)，则
\[
\cos\varphi=\frac{|\overrightarrow{EC_1}\cdot\overrightarrow{FD_1}|}{|\overrightarrow{EC_1}||\overrightarrow{FD_1}|}=\frac{6}{\sqrt{14}\sqrt{24}}=\frac{\sqrt{21}}{14}
\]
\end{enumerate}
\end{solution}

\begin{problem}
设函数 \( f(x) = \left|1 - \frac{1}{x}\right| \)（\(x > 0\)）.

\begin{enumerate}
\item 证明：当 \(0 < a < b\)，且 \(f(a) = f(b)\) 时，\(ab > 1\)；
\item 点 \(P(x_0, y_0)\)（\(0 < x_0 < 1\)）在曲线 \(y = f(x)\) 上，求曲线在点 \(P\) 处的切线与 \(x\) 轴和 \(y\) 轴的正向所围成的三角形面积表达式（用 \(x_0\) 表示）.
\end{enumerate}
\end{problem}

\begin{solution}
\begin{enumerate}
\item 当 \(0<x<1\) 时，\(f(x)=\dfrac{1}{x}-1\)，在 \((0,1)\) 上为减函数；当 \(x>1\) 时，\(f(x)=1-\dfrac{1}{x}\)，在 \((1,+\infty)\) 上为增函数. 又 \(f(1)=0\)，而 \(x\ne1\) 时 \(f(x)>0\)，所以 \(a\)，\(b\) 均不等于 \(1\). 因为 \(0<a<b\) 且 \(f(a)=f(b)\)，\(a\)，\(b\) 不可能在同一个区间 \((0,1)\) 或 \((1,+\infty)\) 内，所以 \(0<a<1<b\). 于是
\(\frac{1}{a}-1=1-\frac{1}{b}\)，从而 \(\frac{1}{a}+\frac{1}{b}=2\)，故 \(ab=\frac{a^2}{2a-1}\).
由 \(b>0\) 得 \(2a-1>0\)，且 \(a\ne1\)，故
\[
ab>1\Longleftrightarrow a^2>2a-1\Longleftrightarrow (a-1)^2>0
\]
所以 \(ab>1\).

\item 因为 \(0<x_0<1\)，所以 \(y_0=f(x_0)=\dfrac{1}{x_0}-1\)，且 \(f'(x_0)=-\dfrac{1}{x_0^2}\). 曲线在 \(P\) 处的切线为 \(y-y_0=-\frac{1}{x_0^2}(x-x_0)\)，整理得 \(y=-\frac{x}{x_0^2}+\frac{2}{x_0}-1\).
令 \(y=0\)，得切线与 \(x\) 轴正向的交点横坐标为 \(x_0(2-x_0)>0\)；令 \(x=0\)，得切线与 \(y\) 轴正向的交点纵坐标为 \(\dfrac{2-x_0}{x_0}>0\). 因而所围三角形面积为
\[
S=\frac12 x_0(2-x_0)\cdot\frac{2-x_0}{x_0}=\frac{(2-x_0)^2}{2}
\]
\end{enumerate}
\end{solution}

\begin{problem}
某中心接到其正东，正西，正北方向三个观测点的报告：正西，正北两个观测点同时听到了一声巨响，正东观测点听到的时间比其他两观测点晚 \(4\mathrm{~s}\). 已知各观测点到该中心的距离都是 \(1020\mathrm{~m}\). 试确定该巨响发生的位置. （假定当时声音传播的速度为 \(340\mathrm{~m/s}\)，相关各点均在同一平面上）
\end{problem}

\begin{solution}
以接报中心为原点 \(O\)，取正东、正北方向分别为 \(x\) 轴、\(y\) 轴正向. 设正西、正东、正北三个观测点分别为 \(A\)，\(B\)，\(C\)，则
\(A=(-1020,0)\)，\(B=(1020,0)\)，\(C=(0,1020)\).
设巨响发生点为 \(P=(x,y)\). 正西、正北观测点同时听到巨响，所以 \(PA=PC\)，从而
\[
(x+1020)^2+y^2=x^2+(y-1020)^2
\]
所以 \(y=-x\).
正东观测点比正西观测点晚 \(4\mathrm{~s}\) 听到声音，而正西观测点与正北观测点同时听到，所以
\[
PB-PA=340\mathrm{~m/s}\cdot4\mathrm{~s}=1360\mathrm{~m}
\]
因此 \(P\) 在以 \(A\)，\(B\) 为焦点的双曲线上. 该双曲线的半焦距为 \(c=1020\)，由焦点距离差得 \(2a=1360\)，即 \(a=680\)，所以
\[
b^2=c^2-a^2=1020^2-680^2=578000
\]
所以双曲线方程为 \(\frac{x^2}{680^2}-\frac{y^2}{578000}=1\). 由于 \(PB>PA\)，点 \(P\) 位于双曲线的左支.
代入 \(y=-x\)，并利用 \(578000=\dfrac54\cdot680^2\)，得
\[
\frac{x^2}{680^2}-\frac{x^2}{578000}=1
\]
化简得 \(x^2=5\cdot680^2\).
故 \(x<0\)，从而 \(y=-x>0\). 所以
\[
P=(-680\sqrt5,680\sqrt5)
\]
该点到中心的距离为 \(OP=680\sqrt{10}\mathrm{~m}\)，方向为接报中心的西偏北 \(45^{\circ}\).
\end{solution}

\begin{problem}
设函数 \( f(x) = x - \ln (x + m) \)，其中常数 \(m\) 为整数.

\begin{enumerate}
\item 当 \(m\) 为何值时，\(f(x) \geqslant 0\)；
\item 定理：若函数 \(g(x)\) 在 \([a, b]\) 上连续，且 \(g(a)\) 与 \(g(b)\) 异号，则至少存在一点 \(x_0 \in (a, b)\)，使 \(g(x_0) = 0\). 试用上述定理证明：当整数 \(m > 1\) 时，方程 \(f(x) = 0\) 在 \([\e^{-m} - m, \e^{2m} - m]\) 内有两个实根.
\end{enumerate}
\end{problem}

\begin{solution}
\begin{enumerate}
\item 函数定义域为 \(x>-m\)，且
\[
f'(x)=1-\frac{1}{x+m}=\frac{x+m-1}{x+m}
\]
因为 \(x+m>0\)，所以 \(f'(x)<0\) 在 \((-m,1-m)\) 上成立，\(f'(x)>0\) 在 \((1-m,+\infty)\) 上成立. 因而 \(f(x)\) 在 \(x=1-m\) 处取得最小值
\[
f(1-m)=1-m
\]
所以对定义域内一切 \(x\) 都有 \(f(x)\geqslant0\) 当且仅当 \(1-m\geqslant0\)，即 \(m\leqslant1\).

\item 设 \(x_1=\e^{-m}-m\)，\(x_0=1-m\)，\(x_2=\e^{2m}-m\). 当整数 \(m>1\) 时，\(\e^{-m}<1<\e^{2m}\)，所以 \(x_1<x_0<x_2\)，且三点均满足 \(x+m>0\). 函数在 \([x_1,x_2]\) 上连续，并且
\(f(x_1)=\e^{-m}>0\)，\(f(x_0)=1-m<0\).
由题给定理，方程 \(f(x)=0\) 在 \((x_1,x_0)\) 内至少有一个实根.

又因为对 \(t>0\) 有 \(\e^t>1+t+\dfrac{t^2}{2}\)，取 \(t=2m\)，所以
\[
f(x_2)=\e^{2m}-3m>1+2m+2m^2-3m=2m^2-m+1>0
\]
再次由题给定理，方程 \(f(x)=0\) 在 \((x_0,x_2)\) 内至少有一个实根. 这两个区间不相交，故所得两个实根不同. 又由第一问中的单调性，\(f\) 在 \((x_1,x_0)\) 和 \((x_0,x_2)\) 上分别严格递减、严格递增，所以每个区间至多有一个实根. 因此方程在 \([\e^{-m}-m,\e^{2m}-m]\) 内恰有两个实根.
\end{enumerate}
\end{solution}

\begin{problem}
设直线 \(l\) 与椭圆 \(\frac{x^2}{25} + \frac{y^2}{16} = 1\) 相交于 \(A\)，\(B\) 两点，\(l\) 又与双曲线 \(x^{2} - y^{2} = 1\) 相交于 \(C\)，\(D\) 两点，\(C\)，\(D\) 三等分线段 \(AB\). 求直线 \(l\) 的方程.
\end{problem}

\begin{solution}
先讨论直线不垂直于 \(x\) 轴的情形，设其方程为 \(y=kx+b\). 将它分别代入椭圆和双曲线，得到关于 \(x\) 的方程
\[
(16+25k^2)x^2+50kbx+25b^2-400=0
\]
\[
(1-k^2)x^2-2kbx-(b^2+1)=0
\]
由于直线与双曲线有两个交点，故 \(k^2\ne1\). 设两条曲线交点的横坐标分别按从小到大排列，则由 \(C\)，\(D\) 三等分线段 \(AB\)，有
\[
x_C+x_D=x_A+x_B
\]
由韦达定理，
\[
-\frac{50kb}{16+25k^2}=\frac{2kb}{1-k^2}
\]
整理得
\[
kb\left(\frac{2}{1-k^2}+\frac{50}{16+25k^2}\right)=0
\]
括号内的分式和为
\[
\frac{2(16+25k^2)+50(1-k^2)}{(1-k^2)(16+25k^2)}=\frac{82}{(1-k^2)(16+25k^2)}\ne0
\]
所以 \(k=0\) 或 \(b=0\).

当 \(k=0\) 时，直线为 \(y=b\). 椭圆和双曲线在该直线上的交点横坐标的绝对值分别为
\(X=5\sqrt{1-\frac{b^2}{16}}\)，\(U=\sqrt{1+b^2}\).
三等分条件给出 \(U=\dfrac{X}{3}\)，故
\[
1+b^2=\frac{25}{9}\left(1-\frac{b^2}{16}\right)\Longrightarrow b^2=\frac{256}{169}
\]
因此得到直线 \(y=\dfrac{16}{13}\) 和 \(y=-\dfrac{16}{13}\).

当 \(b=0\) 时，直线为 \(y=kx\). 此时双曲线方程化为 \((1-k^2)x^2=1\)，要有两个交点必须满足 \(|k|<1\). 椭圆和双曲线交点横坐标的绝对值分别为 \(X=\frac{20}{\sqrt{16+25k^2}}\)，\(U=\frac{1}{\sqrt{1-k^2}}\).
三等分条件给出 \(U=\dfrac{X}{3}\)，故
\[
\frac{1}{1-k^2}=\frac{400}{9(16+25k^2)}\Longrightarrow k^2=\frac{256}{625}
\]
因此得到直线 \(y=\dfrac{16}{25}x\) 和 \(y=-\dfrac{16}{25}x\).

再讨论直线垂直于 \(x\) 轴的情形，设其方程为 \(x=c\). 椭圆和双曲线交点纵坐标的绝对值分别为
\(Y=4\sqrt{1-\frac{c^2}{25}}\)，\(V=\sqrt{c^2-1}\).
三等分条件给出 \(V=\dfrac{Y}{3}\)，故
\[
c^2-1=\frac{16}{9}\left(1-\frac{c^2}{25}\right)\Longrightarrow c^2=\frac{625}{241}
\]
因此得到直线 \(x=\dfrac{25}{\sqrt{241}}\) 和 \(x=-\dfrac{25}{\sqrt{241}}\).

综上，直线 \(l\) 的方程为
\(y=\pm\frac{16}{13}\)，\(y=\pm\frac{16}{25}x\) 或 \(x=\pm\frac{25}{\sqrt{241}}\).
\end{solution}
